\documentclass[11pt,a4paper]{article}

% ============================================================
% PACOTES
% ============================================================

\usepackage[utf8]{inputenc}
\usepackage[T1]{fontenc}
\usepackage[portuguese]{babel}
\usepackage{lmodern}

\usepackage[
a4paper,
top=1.8cm,
bottom=1.8cm,
left=1.8cm,
right=1.8cm
]{geometry}

\usepackage{amsmath,amssymb}
\usepackage{array}
\usepackage{booktabs}
\usepackage{tabularx}
\usepackage{multirow}
\usepackage{multicol}
\usepackage{enumitem}
\usepackage{float}
\usepackage{xcolor}
\usepackage{graphicx}

% ------------------------------------------------------------
% SIUNITX - conforme solicitado
% ------------------------------------------------------------
\usepackage{siunitx}
\sisetup{output-decimal-marker={,}}

% Fórmulas químicas
\usepackage[version=4]{mhchem}

% TikZ e gráficos
\usepackage{tikz}
\usetikzlibrary{
	arrows.meta,
	positioning,
	calc,
	shapes.geometric,
	decorations.pathmorphing,
	backgrounds
}

\usepackage{pgfplots}
\pgfplotsset{compat=1.18}

% Caixas
\usepackage[most]{tcolorbox}

% Cabeçalhos
\usepackage{fancyhdr}
\usepackage{lastpage}

\usepackage{hyperref}

% ============================================================
% CORES
% ============================================================

\definecolor{eublue}{HTML}{145A86}
\definecolor{eudark}{HTML}{17324D}
\definecolor{eulightblue}{HTML}{E4F2FA}
\definecolor{eugreen}{HTML}{3F7D58}
\definecolor{eulightgreen}{HTML}{E8F4EC}
\definecolor{euorange}{HTML}{E28A2B}
\definecolor{eulightorange}{HTML}{FFF0DE}
\definecolor{eured}{HTML}{B33A3A}
\definecolor{eupurple}{HTML}{67518A}
\definecolor{eugray}{HTML}{F2F4F6}
\definecolor{euyellow}{HTML}{F2C94C}

% ============================================================
% CONFIGURAÇÕES GERAIS
% ============================================================

\setlength{\parindent}{0pt}
\setlength{\parskip}{4pt}

\renewcommand{\arraystretch}{1.35}

\hypersetup{
	colorlinks=true,
	linkcolor=eublue,
	urlcolor=eublue
}

% ============================================================
% CABEÇALHO E RODAPÉ
% ============================================================

\pagestyle{fancy}
\fancyhf{}

\fancyhead[L]{
	\textbf{\textcolor{eublue}{EUFISICA exercícios}}
}

\fancyhead[R]{
	\textcolor{eudark}{Física e Química A -- 10.º ano}
}

\fancyfoot[L]{Prof. José Gonçalves}
\fancyfoot[C]{EUFISICA}
\fancyfoot[R]{Página \thepage\ de \pageref{LastPage}}

\renewcommand{\headrulewidth}{0.6pt}
\renewcommand{\footrulewidth}{0.4pt}

% ============================================================
% CAIXAS
% ============================================================

\newcounter{exercicio}

\newenvironment{exercicio}[1]
{
	\refstepcounter{exercicio}
	\begin{tcolorbox}[
		enhanced,
		breakable,
		colback=white,
		colframe=eublue,
		boxrule=0.8pt,
		arc=2mm,
		left=3mm,
		right=3mm,
		top=2mm,
		bottom=2mm,
		title={
			\textbf{Exercício \theexercicio}
			\hspace{1mm}--\hspace{1mm}#1
		},
		coltitle=white,
		colbacktitle=eublue,
		fonttitle=\bfseries
		]
	}
	{
	\end{tcolorbox}
}

\newtcolorbox{recorda}{
	enhanced,
	breakable,
	colback=eulightblue,
	colframe=eublue,
	title=\textbf{Recorda},
	coltitle=eudark,
	colbacktitle=eulightblue,
	fonttitle=\bfseries,
	arc=2mm
}

\newtcolorbox{dados}{
	enhanced,
	breakable,
	colback=eugray,
	colframe=eudark,
	title=\textbf{Dados},
	fonttitle=\bfseries,
	arc=2mm
}

\newtcolorbox{stem}{
	enhanced,
	breakable,
	colback=eulightgreen,
	colframe=eugreen,
	title=\textbf{STEM | Ciência, Tecnologia e Mundo Real},
	coltitle=eugreen,
	colbacktitle=eulightgreen,
	fonttitle=\bfseries,
	arc=2mm
}

\newtcolorbox{desafio}{
	enhanced,
	breakable,
	colback=eulightorange,
	colframe=euorange,
	title=\textbf{Desafio},
	coltitle=euorange,
	colbacktitle=eulightorange,
	fonttitle=\bfseries,
	arc=2mm
}

% ============================================================
% LINHAS DE RESPOSTA
% ============================================================

\newcommand{\linha}{
	\vspace{3mm}
	\noindent\makebox[\linewidth]{\dotfill}
}

\newcommand{\linhas}[1]{
	\par
	\foreach \i in {1,...,#1}{
		\vspace{3.5mm}
		\noindent\makebox[\linewidth]{\dotfill}\par
	}
}

% ============================================================
% DOCUMENTO
% ============================================================

\begin{document}
	
	% ============================================================
	% CAPA / CABEÇALHO DA FICHA
	% ============================================================
	
	\begin{tcolorbox}[
		enhanced,
		colback=eudark,
		colframe=eudark,
		arc=3mm,
		boxrule=0pt,
		left=5mm,
		right=5mm,
		top=5mm,
		bottom=5mm
		]
		
		\begin{center}
			
			{\Huge\bfseries\color{white} EUFISICA exercícios}
			
			\vspace{2mm}
			
			{\Large\bfseries\color{white}
				Estrutura eletrónica dos átomos}
			
			\vspace{2mm}
			
			{\large\color{eulightblue}
				Quantidade de matéria • Massa molar • Espetros atómicos}
			
			{\large\color{eulightblue}
				Átomo de hidrogénio • Energia dos eletrões • Distribuição eletrónica}
			
			\vspace{4mm}
			
			{\large\color{white}\textbf{Física e Química A -- 10.º ano}}
			
			\vspace{2mm}
			
			{\color{white}\textbf{Prof. José Gonçalves}}
			
		\end{center}
		
	\end{tcolorbox}
	
	\vspace{3mm}
	
	\begin{tabularx}{\textwidth}{|X|X|X|}
		\hline
		\textbf{Nome:} &
		\textbf{Turma:} &
		\textbf{Data:} \\[6mm]
		\hline
	\end{tabularx}
	
	\vspace{5mm}
	
	% ============================================================
	% OBJETIVOS
	% ============================================================
	
	\begin{tcolorbox}[
		colback=eugray,
		colframe=eudark,
		title=\textbf{Objetivos de aprendizagem}
		]
		
		No final desta ficha deverás ser capaz de:
		
		\begin{itemize}[leftmargin=6mm,itemsep=1mm]
			
			\item relacionar massa, quantidade de matéria e massa molar;
			
			\item determinar o número de entidades presentes numa determinada
			quantidade de matéria;
			
			\item interpretar espetros de emissão e de absorção;
			
			\item relacionar a energia de um fotão com a frequência e o comprimento
			de onda da radiação;
			
			\item interpretar o espetro do átomo de hidrogénio com base em
			transições eletrónicas;
			
			\item determinar a energia de diferentes níveis do átomo de hidrogénio;
			
			\item interpretar energias de remoção eletrónica e espetros
			fotoeletrónicos;
			
			\item escrever distribuições eletrónicas de átomos e iões;
			
			\item aplicar os princípios de Aufbau, Pauli e Hund.
			
		\end{itemize}
		
	\end{tcolorbox}
	
	% ============================================================
	% DADOS
	% ============================================================
	
	\begin{dados}
		
		Sempre que necessário, utiliza:
		
		\[
		N_A=\SI{6.022e23}{\per\mole}
		\]
		
		\[
		h=\SI{6.626e-34}{\joule\second}
		\]
		
		\[
		c=\SI{3.00e8}{\meter\per\second}
		\]
		
		\[
		1\ \mathrm{eV}
		=
		\SI{1.602e-19}{\joule}
		\]
		
		\[
		E_{\text{fotão}}=h\nu
		\]
		
		\[
		c=\lambda\nu
		\]
		
		logo,
		
		\[
		\boxed{
			E_{\text{fotão}}
			=
			\frac{hc}{\lambda}
		}
		\]
		
	\end{dados}
	
	% ============================================================
	% PARTE 1
	% ============================================================
	
	\section*{\textcolor{eublue}{1. Quantidade de matéria e massa molar}}
	
	\begin{recorda}
		
		A \textbf{quantidade de matéria}, \(n\), exprime-se no SI em
		\textbf{mol}.
		
		Um mol contém:
		
		\[
		\boxed{
			N_A=\SI{6.022e23}{entidades\per\mole}
		}
		\]
		
		A relação entre o número de entidades, \(N\), e a quantidade de matéria
		é:
		
		\[
		\boxed{
			N=nN_A
		}
		\]
		
		A relação entre massa, \(m\), quantidade de matéria, \(n\), e massa
		molar, \(M\), é:
		
		\[
		\boxed{
			n=\frac{m}{M}
		}
		\]
		
		ou
		
		\[
		\boxed{
			m=nM
		}
		\]
		
		A massa molar exprime-se normalmente em:
		
		\[
		\si{\gram\per\mole}
		\]
		
	\end{recorda}
	
	% ============================================================
	% DIAGRAMA MOLE
	% ============================================================
	
	\begin{figure}[H]
		\centering
		
		\begin{tikzpicture}[
			box/.style={
				draw,
				rounded corners,
				minimum width=4.0cm,
				minimum height=1.4cm,
				align=center,
				font=\bfseries
			}
			]
			
			\node[
			box,
			fill=eulightblue,
			draw=eublue
			] (massa) at (0,0)
			{Massa\\\(m\)};
			
			\node[
			box,
			fill=eulightgreen,
			draw=eugreen
			] (mol) at (6,0)
			{Quantidade de matéria\\\(n\)};
			
			\node[
			box,
			fill=eulightorange,
			draw=euorange
			] (part) at (12,0)
			{Número de entidades\\\(N\)};
			
			\draw[
			-{Latex[length=3mm]},
			very thick,
			eublue
			]
			(massa) -- node[above] {\(\div M\)} (mol);
			
			\draw[
			-{Latex[length=3mm]},
			very thick,
			eugreen
			]
			(mol) -- node[above] {\(\times N_A\)} (part);
			
			\draw[
			{Latex[length=3mm]}-,
			very thick,
			eublue
			]
			(massa) -- node[below] {\(\times M\)} (mol);
			
			\draw[
			{Latex[length=3mm]}-,
			very thick,
			euorange
			]
			(mol) -- node[below] {\(\div N_A\)} (part);
			
		\end{tikzpicture}
		
		\caption{Relações entre massa, quantidade de matéria e número de entidades.}
	\end{figure}
	
	% ============================================================
	
	\begin{exercicio}{Quantidade de matéria}
		
		Completa a tabela.
		
		\begin{center}
			
			\begin{tabular}{|c|c|c|}
				\hline
				
				\textbf{Quantidade de matéria}
				&
				\textbf{Número de entidades}
				&
				\textbf{Tipo de entidade}
				\\
				\hline
				
				\(\SI{1.00}{mol}\)
				&
				&
				átomos
				\\[4mm]
				\hline
				
				\(\SI{0.500}{mol}\)
				&
				&
				moléculas
				\\[4mm]
				\hline
				
				\(\SI{2.00}{mol}\)
				&
				&
				iões
				\\[4mm]
				\hline
				
				&
				\(3.011\times10^{23}\)
				&
				moléculas
				\\[4mm]
				\hline
				
				&
				\(1.2044\times10^{24}\)
				&
				átomos
				\\[4mm]
				\hline
				
			\end{tabular}
			
		\end{center}
		
	\end{exercicio}
	
	% ============================================================
	
	\begin{exercicio}{Massa molar}
		
		Considera as seguintes massas atómicas relativas:
		
		\begin{center}
			
			\begin{tabular}{cccccccccc}
				\toprule
				H & C & N & O & Na & Mg & Al & S & Cl & Ca\\
				\midrule
				1,008 & 12,01 & 14,01 & 16,00 & 22,99 &
				24,31 & 26,98 & 32,06 & 35,45 & 40,08\\
				\bottomrule
			\end{tabular}
			
		\end{center}
		
		Determina a massa molar de:
		
		\begin{multicols}{2}
			
			\begin{enumerate}[label=\alph*)]
				
				\item \(\ce{H2O}\)
				
				\item \(\ce{CO2}\)
				
				\item \(\ce{NH3}\)
				
				\item \(\ce{NaCl}\)
				
				\item \(\ce{CaCO3}\)
				
				\item \(\ce{Al2(SO4)3}\)
				
			\end{enumerate}
			
		\end{multicols}
		
		\linhas{5}
		
	\end{exercicio}
	
	% ============================================================
	
	\begin{exercicio}{Da massa ao número de moléculas}
		
		Uma amostra contém \(\SI{18.0}{g}\) de água.
		
		\begin{enumerate}[label=\alph*)]
			
			\item Determina a massa molar da água.
			
			\linhas{1}
			
			\item Determina a quantidade de matéria existente na amostra.
			
			\linhas{2}
			
			\item Determina o número de moléculas de água.
			
			\linhas{2}
			
			\item Determina o número total de átomos de hidrogénio presentes
			na amostra.
			
			\linhas{2}
			
		\end{enumerate}
		
	\end{exercicio}
	
	% ============================================================
	
	\begin{exercicio}{Um composto do quotidiano}
		
		A cafeína apresenta fórmula molecular:
		
		\[
		\ce{C8H10N4O2}
		\]
		
		\begin{enumerate}[label=\alph*)]
			
			\item Determina a massa molar da cafeína.
			
			\linhas{2}
			
			\item Uma bebida contém \(\SI{80}{mg}\) de cafeína.
			Determina a quantidade de matéria de cafeína ingerida.
			
			\linhas{2}
			
			\item Determina o número aproximado de moléculas de cafeína
			correspondente.
			
			\linhas{2}
			
		\end{enumerate}
		
	\end{exercicio}
	
	% ============================================================
	% PARTE 2
	% ============================================================
	
	\newpage
	
	\section*{\textcolor{eublue}{2. Espetros de emissão e de absorção}}
	
	\begin{recorda}
		
		Os átomos possuem níveis de energia eletrónica bem definidos.
		
		Quando um eletrão passa de um nível de energia mais elevado para um
		nível de energia inferior, pode ocorrer a \textbf{emissão de um fotão}.
		
		Quando um átomo absorve um fotão com energia adequada, um eletrão pode
		passar para um nível de energia superior.
		
		Assim:
		
		\begin{center}
			\begin{tabular}{lll}
				\toprule
				& \textbf{Transição eletrónica} & \textbf{Fenómeno}\\
				\midrule
				Emissão &
				\(E_{\text{superior}}\rightarrow E_{\text{inferior}}\)
				&
				Emissão de fotão
				\\
				
				Absorção &
				\(E_{\text{inferior}}\rightarrow E_{\text{superior}}\)
				&
				Absorção de fotão
				\\
				\bottomrule
			\end{tabular}
		\end{center}
		
		Cada elemento possui um conjunto característico de linhas espetrais,
		funcionando como uma espécie de \textbf{impressão digital atómica}.
		
	\end{recorda}
	
	% ============================================================
	% ESPECTRO EMISSÃO H
	% ============================================================
	
	\begin{figure}[H]
		\centering
		
		\begin{tikzpicture}[x=0.020cm,y=1cm]
			
			% Fundo
			\fill[black]
			(380,0) rectangle (700,1.3);
			
			% Linhas hidrogénio
			\draw[violet,very thick]
			(410.2,0)--(410.2,1.3);
			
			\draw[blue!70!violet,very thick]
			(434.0,0)--(434.0,1.3);
			
			\draw[cyan!80!blue,very thick]
			(486.1,0)--(486.1,1.3);
			
			\draw[red,very thick]
			(656.3,0)--(656.3,1.3);
			
			% escala
			\draw[->,thick]
			(380,-0.25)--(710,-0.25);
			
			\foreach \x in {400,450,500,550,600,650,700}{
				\draw (\x,-0.32)--(\x,-0.18);
				\node[below,font=\scriptsize] at (\x,-0.32) {\x};
			}
			
			\node at (540,-0.8)
			{\small Comprimento de onda / nm};
			
		\end{tikzpicture}
		
		\caption{Representação simplificada de algumas riscas de emissão do
			hidrogénio na região visível.}
	\end{figure}
	
	% ============================================================
	% ABSORÇÃO
	% ============================================================
	
	\begin{figure}[H]
		\centering
		
		\begin{tikzpicture}[x=0.020cm,y=1cm]
			
			\shade[
			left color=violet!60,
			right color=red!60
			]
			(380,0) rectangle (700,1.3);
			
			\draw[black,very thick]
			(410.2,0)--(410.2,1.3);
			
			\draw[black,very thick]
			(434.0,0)--(434.0,1.3);
			
			\draw[black,very thick]
			(486.1,0)--(486.1,1.3);
			
			\draw[black,very thick]
			(656.3,0)--(656.3,1.3);
			
			\draw[->,thick]
			(380,-0.25)--(710,-0.25);
			
			\foreach \x in {400,450,500,550,600,650,700}{
				\draw (\x,-0.32)--(\x,-0.18);
				\node[below,font=\scriptsize] at (\x,-0.32) {\x};
			}
			
			\node at (540,-0.8)
			{\small Comprimento de onda / nm};
			
		\end{tikzpicture}
		
		\caption{Representação simplificada de um espetro de absorção.}
	\end{figure}
	
	% ============================================================
	
	\begin{exercicio}{Emissão ou absorção?}
		
		Classifica cada situação como relacionada predominantemente com
		\textbf{emissão} ou \textbf{absorção} de radiação.
		
		\begin{enumerate}[label=\alph*)]
			
			\item Um eletrão passa de um nível de energia \(E_4\) para \(E_2\).
			
			\linha
			
			\item Um eletrão passa de \(E_1\) para \(E_3\).
			
			\linha
			
			\item Observa-se uma linha luminosa sobre um fundo escuro.
			
			\linha
			
			\item Observa-se uma linha escura num espetro contínuo.
			
			\linha
			
			\item Um átomo recebe um fotão cuja energia coincide com a diferença
			entre dois níveis eletrónicos.
			
			\linha
			
		\end{enumerate}
		
	\end{exercicio}
	
	% ============================================================
	% COMPARAÇÃO ESPETROS
	% ============================================================
	
	\begin{exercicio}{Identificação de elementos por espetroscopia}
		
		Os seguintes espetros são esquemáticos.
		
		\begin{center}
			
			\begin{tikzpicture}[x=0.016cm,y=1cm]
				
				% H
				\node[anchor=east,font=\bfseries] at (370,3.6) {H};
				
				\fill[black] (380,3.2) rectangle (700,4.0);
				
				\foreach \x in {410,434,486,656}{
					\draw[eublue!40!cyan,very thick]
					(\x,3.2)--(\x,4.0);
				}
				
				% Na
				\node[anchor=east,font=\bfseries] at (370,2.4) {Na};
				
				\fill[black] (380,2.0) rectangle (700,2.8);
				
				\foreach \x in {568,589}{
					\draw[euyellow,very thick]
					(\x,2.0)--(\x,2.8);
				}
				
				% He
				\node[anchor=east,font=\bfseries] at (370,1.2) {He};
				
				\fill[black] (380,0.8) rectangle (700,1.6);
				
				\foreach \x in {447,502,588,668}{
					\draw[euorange,very thick]
					(\x,0.8)--(\x,1.6);
				}
				
				% Amostra
				\node[anchor=east,font=\bfseries] at (370,0) {X};
				
				\fill[black] (380,-0.4) rectangle (700,0.4);
				
				\foreach \x in {410,434,486,568,589,656}{
					\draw[eugreen!60!yellow,very thick]
					(\x,-0.4)--(\x,0.4);
				}
				
				% escala
				\draw[->]
				(380,-0.8)--(710,-0.8);
				
				\foreach \x in {400,450,500,550,600,650,700}{
					\draw (\x,-0.87)--(\x,-0.73);
					\node[below,font=\scriptsize] at (\x,-0.87) {\x};
				}
				
				\node at (540,-1.35)
				{\small \(\lambda/\mathrm{nm}\)};
				
			\end{tikzpicture}
			
		\end{center}
		
		\begin{enumerate}[label=\alph*)]
			
			\item Que elementos químicos parecem estar presentes na amostra X?
			
			\linhas{2}
			
			\item Explica por que motivo as riscas espetrais permitem identificar
			elementos químicos.
			
			\linhas{2}
			
			\item Poderia um espetro de emissão de dois elementos diferentes ser
			completamente igual? Justifica.
			
			\linhas{2}
			
		\end{enumerate}
		
	\end{exercicio}
	
	% ============================================================
	
	\begin{stem}
		
		A análise dos espetros de absorção da luz proveniente das estrelas permite
		identificar os elementos presentes nas suas atmosferas.
		
		Não é necessário recolher diretamente uma amostra da estrela:
		a própria luz transporta informação sobre a matéria que a emitiu
		ou atravessou.
		
		Esta técnica é fundamental em \textbf{Astronomia},
		\textbf{Astrofísica} e \textbf{Química Analítica}.
		
	\end{stem}
	
	% ============================================================
	% PARTE 3
	% ============================================================
	
	\section*{\textcolor{eublue}{3. Energia dos fotões}}
	
	\begin{recorda}
		
		A energia de um fotão relaciona-se com a frequência pela expressão:
		
		\[
		\boxed{
			E=h\nu
		}
		\]
		
		Como:
		
		\[
		c=\lambda\nu
		\]
		
		então:
		
		\[
		\boxed{
			E=\frac{hc}{\lambda}
		}
		\]
		
		Consequentemente, quanto \textbf{menor} for o comprimento de onda,
		\(\lambda\), \textbf{maior} é a energia do fotão.
		
	\end{recorda}
	
	% ============================================================
	
	\begin{exercicio}{Energia de um fotão}
		
		Considera uma radiação de comprimento de onda:
		
		\[
		\lambda=\SI{500}{nm}
		\]
		
		\begin{enumerate}[label=\alph*)]
			
			\item Converte o comprimento de onda para metros.
			
			\linhas{1}
			
			\item Determina a frequência da radiação.
			
			\linhas{2}
			
			\item Determina a energia de um fotão dessa radiação.
			
			\linhas{2}
			
			\item Determina a energia correspondente a um mol desses fotões,
			em \(\si{\kilo\joule\per\mole}\).
			
			\linhas{3}
			
		\end{enumerate}
		
	\end{exercicio}
	
	% ============================================================
	
	\begin{exercicio}{Comparar fotões}
		
		Considera os fotões A, B e C.
		
		\begin{center}
			
			\begin{tabular}{ccc}
				\toprule
				\textbf{Fotão}
				&
				\(\boldsymbol{\lambda}\)
				&
				\textbf{Região aproximada}
				\\
				\midrule
				
				A & \(\SI{400}{nm}\) & violeta\\
				
				B & \(\SI{550}{nm}\) & verde\\
				
				C & \(\SI{650}{nm}\) & vermelha\\
				
				\bottomrule
			\end{tabular}
			
		\end{center}
		
		\begin{enumerate}[label=\alph*)]
			
			\item Qual apresenta maior frequência?
			
			\linha
			
			\item Qual apresenta menor energia?
			
			\linha
			
			\item Ordena os fotões por ordem crescente de energia.
			
			\linha
			
			\item Justifica a tua resposta recorrendo a uma expressão matemática.
			
			\linhas{2}
			
		\end{enumerate}
		
	\end{exercicio}
	
	% ============================================================
	% HIDROGÉNIO
	% ============================================================
	
	\newpage
	
	\section*{\textcolor{eublue}{4. Espetro e níveis de energia do átomo de hidrogénio}}
	
	\begin{recorda}
		
		No modelo quântico mais simples do átomo de hidrogénio, a energia
		associada ao nível de número quântico principal \(n\) pode ser determinada
		por:
		
		\[
		\boxed{
			E_n=
			-\frac{2,18\times10^{-18}}{n^2}
			\quad \mathrm{J}
		}
		\]
		
		ou, utilizando eletrão-volt:
		
		\[
		\boxed{
			E_n=
			-\frac{13,6}{n^2}
			\quad \mathrm{eV}
		}
		\]
		
		onde:
		
		\[
		n=1,2,3,\ldots
		\]
		
		O sinal negativo significa que o eletrão se encontra ligado ao núcleo.
		
		Quando:
		
		\[
		n\rightarrow\infty
		\]
		
		tem-se:
		
		\[
		E=0
		\]
		
		correspondendo à separação completa entre o eletrão e o núcleo.
		
		Para uma transição:
		
		\[
		\boxed{
			\Delta E=E_f-E_i
		}
		\]
		
		onde \(E_i\) é a energia inicial e \(E_f\) a energia final.
		
		Para uma \textbf{absorção}:
		
		\[
		\Delta E>0
		\]
		
		Para uma \textbf{emissão}:
		
		\[
		\Delta E<0
		\]
		
		e a energia do fotão emitido é:
		
		\[
		\boxed{
			E_{\text{fotão}}=
			|\Delta E|
		}
		\]
		
	\end{recorda}
	
	% ============================================================
	% DIAGRAMA DE NÍVEIS
	% ============================================================
	
	\begin{figure}[H]
		\centering
		
		\begin{tikzpicture}[x=1cm,y=0.55cm]
			
			% eixo
			\draw[-{Latex[length=3mm]},thick]
			(0,-14.5)--(0,1.2);
			
			\node[rotate=90,font=\bfseries]
			at (-0.7,-6.5)
			{Energia / eV};
			
			% níveis
			\draw[very thick,eudark]
			(1,-13.6)--(7,-13.6);
			\node[left] at (1,-13.6) {\(n=1\)};
			\node[right] at (7,-13.6) {\(-13,6\)};
			
			\draw[very thick,eudark]
			(1,-3.40)--(7,-3.40);
			\node[left] at (1,-3.40) {\(n=2\)};
			\node[right] at (7,-3.40) {\(-3,40\)};
			
			\draw[very thick,eudark]
			(1,-1.51)--(7,-1.51);
			\node[left] at (1,-1.51) {\(n=3\)};
			\node[right] at (7,-1.51) {\(-1,51\)};
			
			\draw[very thick,eudark]
			(1,-0.85)--(7,-0.85);
			\node[left] at (1,-0.85) {\(n=4\)};
			\node[right] at (7,-0.85) {\(-0,85\)};
			
			\draw[very thick,eudark]
			(1,-0.54)--(7,-0.54);
			\node[left] at (1,-0.54) {\(n=5\)};
			\node[right] at (7,-0.54) {\(-0,54\)};
			
			\draw[dashed,very thick,eured]
			(1,0)--(7,0);
			
			\node[left] at (1,0)
			{\(n=\infty\)};
			
			\node[right,text=eured]
			at (7,0)
			{\(0\)};
			
			% transições
			\draw[
			-{Latex[length=3mm]},
			very thick,
			eublue
			]
			(2.5,-0.85)--(2.5,-3.40);
			
			\node[
			right,
			text=eublue
			]
			at (2.55,-2.1)
			{\small emissão};
			
			\draw[
			-{Latex[length=3mm]},
			very thick,
			eugreen
			]
			(5,-13.6)--(5,-3.40);
			
			\node[
			right,
			text=eugreen
			]
			at (5.05,-8.5)
			{\small absorção};
			
		\end{tikzpicture}
		
		\caption{Níveis de energia do átomo de hidrogénio e exemplos de
			transições eletrónicas.}
	\end{figure}
	
	% ============================================================
	
	\begin{exercicio}{Energia dos níveis do hidrogénio}
		
		Utiliza:
		
		\[
		E_n=
		-\frac{2,18\times10^{-18}}{n^2}\ \mathrm{J}
		\]
		
		Determina a energia correspondente a:
		
		\begin{enumerate}[label=\alph*)]
			
			\item \(n=1\)
			
			\linhas{2}
			
			\item \(n=2\)
			
			\linhas{2}
			
			\item \(n=3\)
			
			\linhas{2}
			
			\item \(n=4\)
			
			\linhas{2}
			
		\end{enumerate}
		
		Compara os valores obtidos com o diagrama de níveis de energia.
		
	\end{exercicio}
	
	% ============================================================
	
	\begin{exercicio}{Transição eletrónica no átomo de hidrogénio}
		
		Um eletrão do átomo de hidrogénio transita de:
		
		\[
		n=4
		\quad\longrightarrow\quad
		n=2
		\]
		
		\begin{enumerate}[label=\alph*)]
			
			\item Determina \(E_4\).
			
			\linhas{2}
			
			\item Determina \(E_2\).
			
			\linhas{2}
			
			\item Calcula:
			
			\[
			\Delta E=E_f-E_i
			\]
			
			\linhas{2}
			
			\item A transição corresponde a absorção ou emissão?
			
			\linha
			
			\item Determina a energia do fotão envolvido.
			
			\linhas{2}
			
			\item Determina a frequência do fotão.
			
			\linhas{2}
			
			\item Determina o comprimento de onda da radiação emitida.
			
			\linhas{3}
			
			\item Em que região do espetro eletromagnético se encontra esta
			radiação?
			
			\linha
			
		\end{enumerate}
		
	\end{exercicio}
	
	% ============================================================
	
	\begin{exercicio}{Outra risca do hidrogénio}
		
		Um eletrão do átomo de hidrogénio passa de \(n=3\) para \(n=2\).
		
		\begin{enumerate}[label=\alph*)]
			
			\item Representa a transição no diagrama de níveis de energia.
			
			\item Calcula a variação de energia do átomo.
			
			\linhas{2}
			
			\item Determina a energia do fotão emitido.
			
			\linhas{2}
			
			\item Determina o comprimento de onda da radiação emitida.
			
			\linhas{3}
			
			\item Compara este valor com uma das riscas apresentadas anteriormente
			no espetro de emissão do hidrogénio.
			
			\linhas{2}
			
		\end{enumerate}
		
	\end{exercicio}
	
	% ============================================================
	
	\begin{exercicio}{Absorção de radiação}
		
		Um átomo de hidrogénio encontra-se no estado fundamental.
		
		Um fotão provoca a transição:
		
		\[
		n=1\rightarrow n=3
		\]
		
		\begin{enumerate}[label=\alph*)]
			
			\item Determina a energia inicial do eletrão.
			
			\linhas{2}
			
			\item Determina a energia final.
			
			\linhas{2}
			
			\item Determina a energia mínima do fotão que poderá provocar esta
			transição.
			
			\linhas{2}
			
			\item Determina a frequência desse fotão.
			
			\linhas{2}
			
			\item Explica por que motivo um fotão com energia inferior não consegue
			provocar esta transição.
			
			\linhas{3}
			
		\end{enumerate}
		
	\end{exercicio}
	
	% ============================================================
	
	\begin{desafio}
		
		Um eletrão encontra-se inicialmente no nível \(n=2\) do átomo de
		hidrogénio.
		
		Determina a energia mínima necessária para remover completamente esse
		eletrão do átomo.
		
		Recorda que, para um eletrão livre:
		
		\[
		n\rightarrow\infty
		\qquad\Rightarrow\qquad
		E=0
		\]
		
		Apresenta o resultado:
		
		\begin{enumerate}[label=\alph*)]
			
			\item em joule por eletrão;
			
			\item em eletrão-volt por eletrão;
			
			\item em \(\si{\kilo\joule\per\mole}\).
			
		\end{enumerate}
		
		\linhas{5}
		
	\end{desafio}
	
	% ============================================================
	% FOTOELETRÓNICO
	% ============================================================
	
	\newpage
	
	\section*{\textcolor{eublue}{5. Energias de remoção eletrónica e espetro fotoeletrónico}}
	
	\begin{recorda}
		
		Na espetroscopia fotoeletrónica, uma radiação de energia conhecida incide
		sobre os átomos e pode provocar a remoção de eletrões.
		
		Por conservação da energia:
		
		\[
		E_{\text{fotão}}
		=
		E_{\text{remoção}}
		+
		E_{\text{cinética}}
		\]
		
		Assim:
		
		\[
		\boxed{
			E_{\text{remoção}}
			=
			E_{\text{fotão}}
			-
			E_{\text{cinética}}
		}
		\]
		
		ou:
		
		\[
		\boxed{
			E_{\text{remoção}}
			=
			h\nu-E_c
		}
		\]
		
		Quanto mais fortemente ligado estiver um eletrão ao núcleo,
		\textbf{maior será a energia necessária para o remover}.
		
		Num espetro fotoeletrónico:
		
		\begin{itemize}
			
			\item a posição dos picos relaciona-se com a energia de remoção;
			
			\item a intensidade ou área relativa dos picos relaciona-se com o número
			de eletrões existentes em cada subnível;
			
			\item eletrões mais próximos do núcleo apresentam, geralmente,
			maior energia de remoção.
			
		\end{itemize}
		
	\end{recorda}
	
	% ============================================================
	
	\begin{exercicio}{Energia de remoção}
		
		Um fotão com energia:
		
		\[
		E_{\text{fotão}}
		=
		\SI{2.40e-17}{J}
		\]
		
		remove um eletrão que adquire energia cinética:
		
		\[
		E_c=
		\SI{7.0e-18}{J}
		\]
		
		\begin{enumerate}[label=\alph*)]
			
			\item Determina a energia de remoção do eletrão.
			
			\linhas{2}
			
			\item Converte esse valor em \(\mathrm{eV}\).
			
			\linhas{2}
			
			\item Determina a energia de remoção correspondente a um mol de
			eletrões, em \(\si{\kilo\joule\per\mole}\).
			
			\linhas{3}
			
		\end{enumerate}
		
	\end{exercicio}
	
	% ============================================================
	% ESPECTRO FOTOELETRÓNICO ESQUEMÁTICO
	% ============================================================
	
	\begin{figure}[H]
		\centering
		
		\begin{tikzpicture}[x=1.5cm,y=0.7cm]
			
			% Eixos
			\draw[-{Latex[length=3mm]},thick]
			(0,0)--(8.3,0);
			
			\draw[-{Latex[length=3mm]},thick]
			(0,0)--(0,7.5);
			
			\node[rotate=90,font=\bfseries]
			at (-0.65,3.7)
			{Intensidade relativa};
			
			\node[font=\bfseries]
			at (4.1,-1.0)
			{Energia de remoção};
			
			\draw[
			-{Latex[length=3mm]},
			thick,
			eured
			]
			(7,-0.55)--(1,-0.55);
			
			\node[
			text=eured,
			font=\small
			]
			at (4,-0.55)
			{energia de remoção crescente};
			
			% picos Na 1s2 2s2 2p6 3s1
			
			% 1s
			\draw[eublue,very thick]
			(1.2,0)--(1.2,2);
			
			\node[above,font=\bfseries]
			at (1.2,2)
			{\(1s\)};
			
			% 2s
			\draw[eublue,very thick]
			(3.0,0)--(3.0,2);
			
			\node[above,font=\bfseries]
			at (3.0,2)
			{\(2s\)};
			
			% 2p
			\draw[eublue,very thick]
			(4.8,0)--(4.8,6);
			
			\node[above,font=\bfseries]
			at (4.8,6)
			{\(2p\)};
			
			% 3s
			\draw[eublue,very thick]
			(7.0,0)--(7.0,1);
			
			\node[above,font=\bfseries]
			at (7.0,1)
			{\(3s\)};
			
		\end{tikzpicture}
		
		\caption{Espetro fotoeletrónico esquemático de um átomo com distribuição
			eletrónica \(1s^2\,2s^2\,2p^6\,3s^1\). As posições não estão à escala.}
	\end{figure}
	
	% ============================================================
	
	\begin{exercicio}{Interpretar um espetro fotoeletrónico}
		
		Analisa o espetro anterior.
		
		\begin{enumerate}[label=\alph*)]
			
			\item Quantos subníveis ocupados são identificáveis?
			
			\linha
			
			\item Qual dos subníveis apresenta eletrões com maior energia de remoção?
			
			\linha
			
			\item Qual apresenta eletrões mais facilmente removíveis?
			
			\linha
			
			\item Por que motivo o pico correspondente ao subnível \(2p\) é maior
			do que o correspondente ao subnível \(2s\)?
			
			\linhas{2}
			
			\item Quantos eletrões existem no átomo representado?
			
			\linha
			
			\item Identifica o elemento químico.
			
			\linha
			
		\end{enumerate}
		
	\end{exercicio}
	
	% ============================================================
	
	\begin{exercicio}{Energias de ionização sucessivas}
		
		Considere os seguintes valores aproximados relativos ao magnésio.
		
		\begin{center}
			
			\begin{tabular}{cc}
				\toprule
				
				\textbf{Eletrão removido}
				&
				\textbf{Energia de ionização / \(\si{\kilo\joule\per\mole}\)}
				\\
				\midrule
				
				1.º & 738\\
				
				2.º & 1451\\
				
				3.º & 7733\\
				
				4.º & 10543\\
				
				5.º & 13630\\
				
				\bottomrule
			\end{tabular}
			
		\end{center}
		
		\begin{enumerate}[label=\alph*)]
			
			\item Entre que remoções se observa o maior aumento de energia?
			
			\linha
			
			\item O que sugere esse aumento relativamente ao número de eletrões de
			valência do magnésio?
			
			\linhas{2}
			
			\item Sabendo que \(Z(\ce{Mg})=12\), escreve a distribuição eletrónica
			do magnésio.
			
			\linhas{2}
			
			\item Justifica a formação habitual do ião \(\ce{Mg^{2+}}\).
			
			\linhas{2}
			
		\end{enumerate}
		
	\end{exercicio}
	
	% ============================================================
	% DISTRIBUIÇÕES
	% ============================================================
	
	\section*{\textcolor{eublue}{6. Distribuição eletrónica}}
	
	\begin{recorda}
		
		O preenchimento dos orbitais obedece a vários princípios.
		
		\textbf{Princípio da energia mínima ou princípio de Aufbau}
		
		Os eletrões ocupam primeiro os orbitais de menor energia.
		
		\vspace{2mm}
		
		\textbf{Princípio de exclusão de Pauli}
		
		Cada orbital pode conter, no máximo, dois eletrões com spins opostos.
		
		\vspace{2mm}
		
		\textbf{Regra de Hund}
		
		Em orbitais de igual energia, os eletrões ocupam primeiro os orbitais
		separadamente, com spins paralelos, antes de ocorrer emparelhamento.
		
		\vspace{2mm}
		
		A ordem de preenchimento utilizada para os primeiros elementos é:
		
		\[
		\boxed{
			1s
			\rightarrow
			2s
			\rightarrow
			2p
			\rightarrow
			3s
			\rightarrow
			3p
			\rightarrow
			4s
		}
		\]
		
	\end{recorda}
	
	% ============================================================
	% DIAGRAMA ORBITAIS
	% ============================================================
	
	\begin{figure}[H]
		\centering
		
		\begin{tikzpicture}[x=1cm,y=1cm]
			
			\node[font=\bfseries] at (-1,4.8) {Energia};
			
			\draw[-{Latex[length=3mm]},thick]
			(-1,0)--(-1,4.5);
			
			% 1s
			\draw[eublue,thick]
			(0,0.4) rectangle (0.8,1.0);
			
			\node[left] at (-0.1,0.7) {\(1s\)};
			
			% 2s
			\draw[eublue,thick]
			(0,1.5) rectangle (0.8,2.1);
			
			\node[left] at (-0.1,1.8) {\(2s\)};
			
			% 2p
			\foreach \x in {0,1,2}{
				\draw[eugreen,thick]
				({1.5+0.9*\x},2.2)
				rectangle
				({2.2+0.9*\x},2.8);
			}
			
			\node[left] at (1.35,2.5) {\(2p\)};
			
			% 3s
			\draw[eublue,thick]
			(0,3.0) rectangle (0.8,3.6);
			
			\node[left] at (-0.1,3.3) {\(3s\)};
			
			% 3p
			\foreach \x in {0,1,2}{
				\draw[eugreen,thick]
				({1.5+0.9*\x},3.7)
				rectangle
				({2.2+0.9*\x},4.3);
			}
			
			\node[left] at (1.35,4.0) {\(3p\)};
			
		\end{tikzpicture}
		
		\caption{Representação esquemática de orbitais \(s\) e \(p\).}
	\end{figure}
	
	% ============================================================
	
	\begin{exercicio}{Escrever distribuições eletrónicas}
		
		Escreve a distribuição eletrónica, no estado fundamental, dos seguintes
		átomos.
		
		\begin{center}
			
			\begin{tabular}{ccc}
				\toprule
				\textbf{Elemento} &
				\(\boldsymbol{Z}\) &
				\textbf{Distribuição eletrónica}
				\\
				\midrule
				
				\ce{C} & 6 & \\[4mm]
				
				\ce{O} & 8 & \\[4mm]
				
				\ce{Ne} & 10 & \\[4mm]
				
				\ce{Na} & 11 & \\[4mm]
				
				\ce{Mg} & 12 & \\[4mm]
				
				\ce{Al} & 13 & \\[4mm]
				
				\ce{Cl} & 17 & \\[4mm]
				
				\ce{Ca} & 20 & \\[4mm]
				
				\bottomrule
			\end{tabular}
			
		\end{center}
		
	\end{exercicio}
	
	% ============================================================
	
	\begin{exercicio}{Diagramas de orbitais}
		
		Representa os eletrões nos diagramas de caixas respeitando os princípios
		de Pauli e Hund.
		
		\subsection*{a) Carbono, \(Z=6\)}
		
		\[
		1s^2\,2s^2\,2p^2
		\]
		
		\begin{center}
			
			\begin{tikzpicture}[scale=0.9]
				
				\node at (-0.8,0.45) {\(1s\)};
				
				\draw (0,0) rectangle (1,0.9);
				
				\node at (1.7,0.45) {\(2s\)};
				
				\draw (2.5,0) rectangle (3.5,0.9);
				
				\node at (4.2,0.45) {\(2p\)};
				
				\foreach \x in {5,6,7}{
					\draw (\x,0) rectangle ({\x+1},0.9);
				}
				
			\end{tikzpicture}
			
		\end{center}
		
		\vspace{5mm}
		
		\subsection*{b) Oxigénio, \(Z=8\)}
		
		\[
		1s^2\,2s^2\,2p^4
		\]
		
		\begin{center}
			
			\begin{tikzpicture}[scale=0.9]
				
				\node at (-0.8,0.45) {\(1s\)};
				
				\draw (0,0) rectangle (1,0.9);
				
				\node at (1.7,0.45) {\(2s\)};
				
				\draw (2.5,0) rectangle (3.5,0.9);
				
				\node at (4.2,0.45) {\(2p\)};
				
				\foreach \x in {5,6,7}{
					\draw (\x,0) rectangle ({\x+1},0.9);
				}
				
			\end{tikzpicture}
			
		\end{center}
		
	\end{exercicio}
	
	% ============================================================
	
	\begin{exercicio}{Detetar erros}
		
		Um aluno escreveu as seguintes distribuições eletrónicas:
		
		\[
		\text{A:}\qquad
		1s^2\,2s^2\,2p^7
		\]
		
		\[
		\text{B:}\qquad
		1s^2\,2s^2\,2p^6\,3s^3
		\]
		
		\[
		\text{C:}\qquad
		1s^2\,2s^2\,2p^6\,3s^2\,3p^5
		\]
		
		\begin{enumerate}[label=\alph*)]
			
			\item Qual ou quais das distribuições são impossíveis?
			
			\linhas{1}
			
			\item Indica a regra violada em cada caso impossível.
			
			\linhas{2}
			
			\item A distribuição C corresponde a que elemento químico?
			
			\linha
			
			\item Quantos eletrões de valência possui esse elemento?
			
			\linha
			
		\end{enumerate}
		
	\end{exercicio}
	
	% ============================================================
	
	\begin{exercicio}{Distribuição eletrónica de iões}
		
		Escreve a distribuição eletrónica das seguintes espécies.
		
		\begin{enumerate}[label=\alph*)]
			
			\item \(\ce{Na+}\)
			
			\[
			Z(\ce{Na})=11
			\]
			
			\linhas{1}
			
			\item \(\ce{Mg^{2+}}\)
			
			\[
			Z(\ce{Mg})=12
			\]
			
			\linhas{1}
			
			\item \(\ce{Cl-}\)
			
			\[
			Z(\ce{Cl})=17
			\]
			
			\linhas{1}
			
			\item \(\ce{O^{2-}}\)
			
			\[
			Z(\ce{O})=8
			\]
			
			\linhas{1}
			
			\item O que têm em comum as quatro distribuições eletrónicas obtidas?
			
			\linhas{2}
			
		\end{enumerate}
		
	\end{exercicio}
	
	% ============================================================
	% QUESTÃO INTEGRADORA
	% ============================================================
	
	\newpage
	
	\section*{\textcolor{euorange}{7. Problema integrador}}
	
	\begin{desafio}
		
		Uma equipa de investigação analisa a luz emitida por uma amostra de
		hidrogénio excitado.
		
		Foi observada uma linha espetral com comprimento de onda próximo de:
		
		\[
		\lambda=\SI{486}{nm}
		\]
		
		Considere ainda que a amostra contém:
		
		\[
		m=\SI{2.00}{mg}
		\]
		
		de átomos de hidrogénio.
		
		Admite:
		
		\[
		M(\ce{H})=\SI{1.008}{g\per\mole}
		\]
		
		Responde às questões seguintes.
		
		\begin{enumerate}[label=\alph*)]
			
			\item Determina a frequência da radiação observada.
			
			\linhas{2}
			
			\item Determina a energia de um fotão dessa radiação.
			
			\linhas{2}
			
			\item No átomo de hidrogénio, a risca corresponde à transição:
			
			\[
			n=4\rightarrow n=2
			\]
			
			Determina \(E_4\).
			
			\linhas{2}
			
			\item Determina \(E_2\).
			
			\linhas{2}
			
			\item Determina a variação de energia do átomo.
			
			\linhas{2}
			
			\item Compara o módulo de \(\Delta E\) com a energia do fotão calculada
			anteriormente.
			
			\linhas{2}
			
			\item Explica a relação física entre os dois valores.
			
			\linhas{3}
			
			\item Determina a quantidade de matéria de hidrogénio existente na
			amostra.
			
			\linhas{2}
			
			\item Determina o número aproximado de átomos existentes na amostra.
			
			\linhas{2}
			
			\item Supondo, apenas como modelo idealizado, que cada átomo emitia um
			fotão correspondente à transição \(n=4\rightarrow n=2\), determina a
			energia total emitida.
			
			\linhas{3}
			
		\end{enumerate}
		
	\end{desafio}
	
	% ============================================================
	% ASTRONOMIA
	% ============================================================
	
	\section*{\textcolor{eugreen}{8. Ciência no mundo real}}
	
	\begin{stem}
		
		\textbf{Como sabemos de que são feitas as estrelas?}
		
		A luz proveniente de uma estrela contém uma enorme quantidade de
		informação.
		
		Ao atravessar as regiões mais externas da estrela, determinados
		comprimentos de onda são absorvidos pelos átomos presentes.
		
		Quando essa luz é analisada num espetroscópio, observam-se linhas escuras
		características.
		
		Comparando essas linhas com os espetros medidos em laboratório, os
		cientistas conseguem identificar elementos presentes nas estrelas,
		mesmo quando se encontram a muitos anos-luz de distância.
		
		A espetroscopia permite ainda estudar:
		
		\begin{itemize}
			
			\item composição química;
			
			\item temperatura;
			
			\item movimento das estrelas;
			
			\item atmosferas planetárias;
			
			\item matéria interestelar.
			
		\end{itemize}
		
	\end{stem}
	
	% ============================================================
	
	\begin{exercicio}{Interpretar o espetro de uma estrela}
		
		Foram identificadas no espetro de absorção de uma estrela linhas próximas
		dos seguintes comprimentos de onda:
		
		\[
		410\,\mathrm{nm},
		\qquad
		434\,\mathrm{nm},
		\qquad
		486\,\mathrm{nm},
		\qquad
		656\,\mathrm{nm}.
		\]
		
		\begin{enumerate}[label=\alph*)]
			
			\item Compara estes valores com o espetro do hidrogénio apresentado
			na ficha.
			
			\item Que conclusão pode ser retirada?
			
			\linhas{2}
			
			\item Explica por que motivo não é necessário recolher fisicamente uma
			amostra da estrela para identificar o hidrogénio.
			
			\linhas{3}
			
		\end{enumerate}
		
	\end{exercicio}
	
	% ============================================================
	% SÍNTESE
	% ============================================================
	
	\section*{\textcolor{eublue}{9. Síntese final}}
	
	Completa o mapa conceptual.
	
	\begin{center}
		
		\begin{tikzpicture}[
			node distance=2.0cm and 2.2cm,
			caixa/.style={
				draw,
				rounded corners,
				minimum width=3.6cm,
				minimum height=1.1cm,
				align=center
			}
			]
			
			\node[
			caixa,
			fill=eulightblue
			] (atom)
			{\textbf{Átomo}};
			
			\node[
			caixa,
			fill=eulightgreen,
			below left=of atom
			] (config)
			{Distribuição\\eletrónica};
			
			\node[
			caixa,
			fill=eulightorange,
			below right=of atom
			] (energy)
			{Níveis de\\energia};
			
			\node[
			caixa,
			fill=eugray,
			below=of config
			] (pes)
			{Espetro\\fotoeletrónico};
			
			\node[
			caixa,
			fill=eugray,
			below=of energy
			] (spectrum)
			{Espetro de emissão\\ou absorção};
			
			\node[
			caixa,
			fill=eulightblue,
			below=of spectrum
			] (photon)
			{Fotões\\\(E=h\nu\)};
			
			\draw[-{Latex},thick]
			(atom)--(config);
			
			\draw[-{Latex},thick]
			(atom)--(energy);
			
			\draw[-{Latex},thick]
			(config)--(pes);
			
			\draw[-{Latex},thick]
			(energy)--(spectrum);
			
			\draw[-{Latex},thick]
			(spectrum)--(photon);
			
		\end{tikzpicture}
		
	\end{center}
	
	% ============================================================
	% AUTOAVALIAÇÃO
	% ============================================================
	
	\section*{\textcolor{eublue}{10. Autoavaliação}}
	
	\begin{center}
		
		\small
		
		\begin{tabularx}{\textwidth}{
				>{\raggedright\arraybackslash}X
				c c c
			}
			
			\toprule
			
			&
			\textbf{Tenho dúvidas}
			&
			\textbf{Estou a aprender}
			&
			\textbf{Consigo sozinho}
			\\
			
			\midrule
			
			Relaciono \(m\), \(n\), \(M\) e \(N\).
			&
			\(\square\)&\(\square\)&\(\square\)
			\\[3mm]
			
			Calculo massas molares.
			&
			\(\square\)&\(\square\)&\(\square\)
			\\[3mm]
			
			Interpreto espetros de emissão e absorção.
			&
			\(\square\)&\(\square\)&\(\square\)
			\\[3mm]
			
			Calculo a energia de um fotão.
			&
			\(\square\)&\(\square\)&\(\square\)
			\\[3mm]
			
			Determino a energia de um nível do átomo de hidrogénio.
			&
			\(\square\)&\(\square\)&\(\square\)
			\\[3mm]
			
			Relaciono \(\Delta E\) com emissão e absorção.
			&
			\(\square\)&\(\square\)&\(\square\)
			\\[3mm]
			
			Interpreto um espetro fotoeletrónico.
			&
			\(\square\)&\(\square\)&\(\square\)
			\\[3mm]
			
			Escrevo distribuições eletrónicas.
			&
			\(\square\)&\(\square\)&\(\square\)
			\\
			
			\bottomrule
			
		\end{tabularx}
		
	\end{center}
	
	\vfill
	
	% ============================================================
	% RODAPÉ VISUAL
	% ============================================================
	
	\begin{center}
		
		\begin{tikzpicture}
			
			\draw[eublue,very thick]
			(0,0)--(13,0);
			
			\node[
			text=eublue,
			font=\Large\bfseries
			]
			at (6.5,0.55)
			{EUFISICA exercícios};
			
			\node[
			text=eudark,
			font=\small
			]
			at (6.5,-0.45)
			{Compreender a Física e a Química para compreender o mundo};
			
			\node[
			text=eudark,
			font=\small
			]
			at (6.5,-0.85)
			{Prof. José Gonçalves};
			
		\end{tikzpicture}
		
	\end{center}
	
\end{document}